\section{Introduction}
A representation can be mathematically valid while conveying almost no distinction between observations. This problem is especially easy to overlook when a model exports a normalized probability vector: every row can sum to one even when all rows assign nearly all their mass to the same component. Topic models define observations through mixtures over latent components~\cite{blei2003latent}, and neural variational constructions provide a practical way to fit related count models~\cite{srivastava2017autoencoding}. Their constraints specify the space of possible outputs, but do not ensure that a trained model uses that space informatively. Before interpreting topic allocations as cell programs, their variation across cells must therefore be examined directly.

Self-attention adds flexible interactions among learned representations, but its scientific interpretation depends on the objects that attend to one another. The transformer architecture introduced attention-based sequence processing~\cite{vaswani2017attention}; an application to expression counts need not use genes, cells, or biological pathways as its tokens. In the model examined here, each cell's full feature vector passes through eight learned projections, producing eight within-cell tokens. Those tokens interact through two attention layers and a pooling query before parameterizing a topic allocation. The model is therefore neither a cross-cell attention system nor a pretrained gene-language model. Stating the tensor axes and projection construction is necessary before asking what a failure says about attention.

The prior regularizer presents another interpretive boundary. A low concentration parameter may encourage sparse topic proportions, but sparsity within each cell is distinct from diversity between cells. Distinct cells can occupy different simplex corners, or occupy the same corner with almost no between-cell information. Perplexity near one cannot distinguish these possibilities; coordinate range, argmax occupancy and native distances supply complementary checks. We call the latter behavior allocation collapse. Posterior collapse in variational inference instead concerns a posterior approaching its prior or failing to use latent variables~\cite{lucas2019collapse}. Constancy of $\operatorname{softmax}(\mu)$ alone does not establish either posterior equality or the behavior of the learned variance.

Distance-based summaries can become misleading near this boundary. An allocation-coherence ratio compares similarity among graph neighbors with similarity among arbitrary cells. When every allocation is nearly identical, both quantities approach zero and their ratio becomes numerically unstable or undefined. A software sentinel returned for this condition is a reporting convention, not a biological measurement. Graph reachability can remain nontrivial because neighbor construction, ties, and graph thresholds operate on a different object. A repeated scalar graph fraction also need not identify the same severed edges. Reliable interpretation requires the status of the native representation, the estimator's denominator, and the graph's identity to remain separate observations.

Low-dimensional visualization creates a further opportunity for confusion. UMAP~\cite{mcinnes2018umap} produces a fitted two-dimensional layout from a representation, and a visually structured layout does not establish that the underlying allocation contains the intended biological information. Label-neighbor purity in that layout has its own definition, including how self-neighbors are excluded. Correcting a neighbor-selection rule can change the displayed score without changing any coordinates or training outcome. Similarly, agreement with Palantir branch probabilities~\cite{setty2019palantir} describes a computational annotation rather than experimentally observed fate. A degenerate allocation should not acquire apparent credibility because a separate visualization or graph summary happens to look plausible.

To determine whether an apparent failure is transient, comparisons must account for optimization exposure. An observation at one short training cap could reflect an unfinished fit, but increasing the cap does not create equal exposure when different runs stop early. Historical endpoints are useful records of actual behavior, provided completed epochs are reported and the same endpoint is not counted repeatedly as a new replication. A new controlled concentration comparison can strengthen this evidence by matching the autoencoder initialization, batch sequence, stochastic draws, and number of optimizer updates. Direct interventions on the trained token tensor can then test whether the implemented pathway responds to token content or inadvertently mixes information between cells.

This study examines an attention-based count topic model across four public expression backgrounds, with particular focus on low-concentration hematopoietic allocations. Historical native arrays permit direct assessment of allocation variation alongside graph and visualization readouts. A new matched concentration experiment in hematopoietic and dentate data tests the sensitivity of native allocations under complete fixed training budgets. Trained-state interventions compare token permutation, token homogenization, token duplication, and changes to the other cells in a batch. These interventions have distinct roles: invariance controls test the computation's intended axes, while content changes test local sensitivity. The objective is a self-contained account of allocation behavior and its measurable boundaries, without inferring that attention architectures generally collapse or that an informative allocation automatically identifies a biological lineage.
